\begin{proposition}[The translation component]
\label{prop:c10-P1}
Under the hypotheses of Definition \ref{def:c10-hyp} with $P_b$ in place of $\Delta$, the component $P_1(O)$ of the master formula of Theorem \ref{thm:c11-classes} (iii) is
\[
P_1(O)=-\big\{\hbar F_1^{(1)}(P_b),O\big\}+\big(\text{composite term of }O\text{ along }P_b\big),
\]
the composite term along $P_b$ defined as along $\Delta$ with the composite kernel $2\big(\partial_b(P(\varepsilon,\infty)\star)^2\otimes1\big)K_0$ (Theorem \ref{thm:c9-P1}), a line from $c$ to $A|_V$ of momentum degree $-2$ without explicit position. There is no term with the coefficient of $\log\varepsilon$, since $\partial_b$ commutes with $P(\varepsilon,\infty)\star$. Its two-body part is the image under $Q$ of the renormalised composite cycles through a pair along $P_b$, local of Weyl weight $w_{X_a}+w_{X_b}+1$. On the part with two legs $A|_{W_{\mathbb R}^\perp}$ of the scalar bilinears $P_1(O)=-(\tfrac13+\theta_P)$ units times $\{S^{(1)}(P_b),O\}$ (Theorem \ref{thm:c9-P1}), which vanishes in the translation-covariant scheme, and on the planar scalar sector Proposition \ref{prop:c10-scalarsector} applies.
\end{proposition}
